\documentclass[11pt,a4paper]{article}
\usepackage{latex/research-report}
\usepackage{newunicodechar}
\newunicodechar{→}{\ensuremath{\rightarrow}}
\newunicodechar{×}{\ensuremath{\times}}
\begin{document}
\hypersetup{pdftitle={Scheduling tasks across different computing resources},pdfauthor={Gleb Mozhaiskii}}
\section{Scheduling tasks across different computing
resources}\label{scheduling-tasks-across-different-computing-resources}

\textbf{Research prototype · UPV/EHU · October 2026}

\textbf{The idea.} An \emph{orchestrator} is a program that receives
computing tasks and decides which available machine should execute each
one. I built a small prototype to study how network delay, computing
power and queues affect that decision.

The results below come from \textbf{three fresh Linux VM sessions,
A01-A03}, each running the same experiment once. All three completed.
They demonstrate a working implementation under repeated starts on one
host; they do not establish which scheduling algorithm is best.

\subsection{Research context and relation to Wang et al.
(2026)}\label{research-context-and-relation-to-wang-et-al.-2026}

Wang et al. {[}1{]} motivate fair comparisons of scheduling algorithms
on the same infrastructure, under different workload conditions, with
service-level outcomes and decision overhead reported separately. I
follow those \textbf{evaluation principles} in a smaller placement
experiment with three CPU-limited workers, actual task queues, and
emulated network delay.

This is \textbf{not a reproduction} of their Continuous Multi-Mode
Scheduling (CMMS) benchmark. Their system chooses both an execution node
and an AI inference mode, and includes output quality in its objectives.
My system chooses only a worker for a fixed CPU task. Their
heterogeneous physical testbed and my single-VM emulation are also not
equivalent. The present results demonstrate a working experimental
prototype, not a general policy advantage.

\subsection{1. The computing
environment}\label{the-computing-environment}

The three worker roles refer to \textbf{where computing resources are
placed in a network}:

\begin{itemize}
\tightlist
\item
  \textbf{Edge:} computing close to where the request originates. It
  usually offers short communication paths, but can have limited
  processing power.
\item
  \textbf{Fog:} an intermediate computing layer between Edge and Cloud.
  Here, it also forwards traffic to the Cloud.
\item
  \textbf{Cloud:} more distant computing resources that can offer higher
  capacity, but may require longer data transfers.
\end{itemize}

These are logical \textbf{roles}, not three separate physical locations
in my experiment.

\textbf{MADT (Distributed Application Modelling System)} is an
open-source tool for deploying applications into an emulated IP-network
topology and controlling network conditions such as delay and bandwidth.
I use \href{https://github.com/dltcspbu/madt}{MADT} to create the
network connections and impairments while running real application code;
this is \textbf{network emulation, not four independent physical sites}.
To reproduce network conditions in a controlled way, I run four isolated
\textbf{Docker containers inside one Linux virtual machine}, connected
through a network emulator. A \emph{CPU quota} limits how much
processing time a container may use.

\begin{figure}[H]
\centering
\includegraphics[width=1\linewidth,height=\textheight,keepaspectratio,alt={System topology}]{build/figures/topology.pdf}
\caption{System topology}
\end{figure}

\begin{table}[H]
\centering
\label{table:topology:start}
\begin{tabular}{p{0.24\linewidth}p{0.54\linewidth}r}
\toprule
\textbf{Container} & \textbf{What it does} & \textbf{CPU quota} \\
\midrule
Gateway / orchestrator & Receives tasks and selects workers & 0.50 CPU \\
Edge worker & Executes tasks locally & 0.25 CPU \\
Fog worker & Executes tasks and forwards Cloud traffic & 0.50 CPU \\
Cloud worker & Executes tasks after traffic passes through Fog & 1.00 CPU \\
\bottomrule
\end{tabular}
\label{table:topology:end}
\end{table}


Each session uses a fresh ARM64 Linux VM with \textbf{4 virtual CPUs, 4
GiB RAM and a 20 GiB sparse disk}. The three VMs ran sequentially on the
same Apple Silicon Mac. They are separate VM sessions, not independent
physical hosts.

Every container is limited to \textbf{256 MiB of memory}. A
\textbf{queue} is where work waits until it can proceed; \emph{FIFO}
means tasks are handled in arrival order (\emph{first in, first out}).
The Gateway has a waiting queue and one outgoing queue per worker; each
worker has a local FIFO and \textbf{one execution slot}. Choosing a
worker does not imply immediate execution.

Traffic to Cloud follows \textbf{Gateway → Fog → Cloud}, while Edge can
be reached directly on the access network.

\subsection{2. What kind of task do I
send?}\label{what-kind-of-task-do-i-send}

I use a \textbf{deterministic CPU-heavy calculation}, not an artificial
sleep.

\textbf{PBKDF2-HMAC-SHA256} is normally a method for deriving a
cryptographic key. It repeatedly applies a SHA-256-based keyed hash for
a specified number of iterations. \textbf{More iterations mean more
repeated CPU work.} Here I use it only as a controlled workload: the
same input and iteration count produce the same output, which lets the
Gateway verify every worker\textquotesingle s answer. I am \emph{not}
evaluating cryptographic security.

Each task combines one of two \textbf{input sizes} (64 KiB or 1 MiB)
with one of two \textbf{CPU work levels} (200,000 or 800,000 hash
iterations). This gives \textbf{four task types}. The input is actually
transferred over the emulated network, then the worker performs the
calculation and returns a verifiable result.

For scoring I use a simple \textbf{one-second deadline}, measured from
the task\textquotesingle s planned arrival. This is an illustrative
experimental target, not a production service requirement.

\subsection{3. How does the orchestrator choose a
worker?}\label{how-does-the-orchestrator-choose-a-worker}

A \emph{scheduling algorithm} decides where to send a task. I
implemented three simple choices:

\begin{itemize}
\tightlist
\item
  \textbf{Round Robin (RR):} takes turns assigning tasks to workers that
  have capacity available. It does not try to predict execution time.
\item
  \textbf{Compute-ECT:} \emph{ECT means Earliest Completion Time}. It
  estimates when each worker would finish the task using its computing
  speed and known queued work, \textbf{without including transfer cost}.
\item
  \textbf{Network-ECT:} uses a similar completion-time estimate, but
  \textbf{also includes expected time to send the input and return the
  result}.
\end{itemize}

Both ECT algorithms make predictions from information available to the
orchestrator at decision time; neither knows future task completion
times. Processing and transfer profiles are calibrated separately before
comparison in each VM, then held fixed. The algorithms and task streams
remain unchanged across sessions, but calibration and measured queue
state can lead to different placements.

I compare these algorithms with two task-arrival patterns:

\begin{itemize}
\tightlist
\item
  \textbf{Steady:} approximately one task every 0.5 seconds.
\item
  \textbf{Burst:} tasks arrive close together in two short bursts,
  creating opportunities for queueing.
\end{itemize}

I compare \textbf{normal networking} with a second condition in which
the network emulator adds a nominal \textbf{20 ms one-way delay to
packets entering Cloud from Fog}. It affects the input path toward
Cloud, not every link or the return path. Measured end-to-end task
latency can change by a different amount because execution and queue
waiting also contribute.

That makes \textbf{3 algorithms × 2 arrival patterns × 2 network
conditions = 12 scenarios}, each offering 12 tasks over 6 seconds.

\subsection{4. Three concrete examples}\label{three-concrete-examples}

Each example holds the arrival pattern and network condition fixed
across policies. \textbf{A01, A02 and A03 are separate VM sessions.}
Each latency is the median of 12 tasks in that session;
\textquotesingle On time\textquotesingle{} gives the descriptive total
out of 36 offers across the three sessions. Those tasks are not 36
independent repetitions.

\Needspace{110pt}

\textbf{Example A - Steady arrivals, normal network.} Tasks are spaced
apart, so the workers have relatively little queueing.

\begin{table}[H]
\centering
\label{table:example-a:start}
\begin{tabular}{lrrrr}
\toprule
\textbf{Scheduler} & \textbf{Run 1 (ms)} & \textbf{Run 2 (ms)} & \textbf{Run 3 (ms)} & \textbf{On time} \\
\midrule
Round Robin & 181.7 & 236.7 & 227.2 & 36/36 \\
Compute-ECT & 118.0 & 121.2 & 119.2 & 36/36 \\
Network-ECT & 113.6 & 120.1 & 119.2 & 36/36 \\
\bottomrule
\end{tabular}
\par\smallskip
{\small Tasks assigned: Edge / Fog / Cloud (12 per run).}\par\smallskip
\begin{tabular}{lccc}
\toprule
\textbf{Scheduler} & \textbf{Run 1} & \textbf{Run 2} & \textbf{Run 3} \\
\midrule
Round Robin & 4 / 4 / 4 & 4 / 4 / 4 & 4 / 4 / 4 \\
Compute-ECT & 0 / 0 / 12 & 0 / 0 / 12 & 0 / 0 / 12 \\
Network-ECT & 0 / 0 / 12 & 0 / 0 / 12 & 0 / 0 / 12 \\
\bottomrule
\end{tabular}
\label{table:example-a:end}
\end{table}


\Needspace{110pt}

\textbf{Example B - Bursty arrivals, normal network.} The arrival stream
is concentrated into two short bursts. The single missed deadline
occurred in A03 with Round Robin.

\begin{table}[H]
\centering
\label{table:example-b:start}
\begin{tabular}{lrrrr}
\toprule
\textbf{Scheduler} & \textbf{Run 1 (ms)} & \textbf{Run 2 (ms)} & \textbf{Run 3 (ms)} & \textbf{On time} \\
\midrule
Round Robin & 185.5 & 171.9 & 189.8 & 35/36 \\
Compute-ECT & 131.2 & 143.3 & 136.0 & 36/36 \\
Network-ECT & 137.3 & 138.3 & 139.0 & 36/36 \\
\bottomrule
\end{tabular}
\par\smallskip
{\small Tasks assigned: Edge / Fog / Cloud (12 per run).}\par\smallskip
\begin{tabular}{lccc}
\toprule
\textbf{Scheduler} & \textbf{Run 1} & \textbf{Run 2} & \textbf{Run 3} \\
\midrule
Round Robin & 4 / 4 / 4 & 4 / 4 / 4 & 4 / 4 / 4 \\
Compute-ECT & 0 / 1 / 11 & 0 / 1 / 11 & 0 / 1 / 11 \\
Network-ECT & 0 / 1 / 11 & 0 / 1 / 11 & 0 / 1 / 11 \\
\bottomrule
\end{tabular}
\label{table:example-b:end}
\end{table}


\Needspace{110pt}

\textbf{Example C - Bursty arrivals, +20 ms toward Cloud.} The input
route to Cloud has additional one-way delay.

\begin{table}[H]
\centering
\label{table:example-c:start}
\begin{tabular}{lrrrr}
\toprule
\textbf{Scheduler} & \textbf{Run 1 (ms)} & \textbf{Run 2 (ms)} & \textbf{Run 3 (ms)} & \textbf{On time} \\
\midrule
Round Robin & 210.6 & 210.4 & 210.2 & 36/36 \\
Compute-ECT & 191.0 & 186.8 & 198.7 & 36/36 \\
Network-ECT & 186.9 & 206.5 & 204.5 & 36/36 \\
\bottomrule
\end{tabular}
\label{table:example-c:end}
\end{table}


\textbf{Interpretation:} Round Robin spreads work across the available
workers. The ECT policies usually choose Cloud, sometimes Fog, under
this small workload. Every observed ECT latency median was below its
matched Round Robin median, but both ECT policies also spent more time
deciding. These are descriptive observations from three sessions, not an
established performance ranking.

\Needspace{260pt}

\subsubsection{Complete set of 12
scenarios}\label{complete-set-of-12-scenarios}

The table includes \textbf{all tested combinations in all three
sessions}. +20 ms means nominal extra one-way Cloud ingress delay. Each
A-column is a 12-task median in milliseconds, measured from planned
arrival to verified response. \textquotesingle On time\textquotesingle{}
retains the one-second deadline and all 36 offers per row. Per-session
placements and individual tasks are available in the public data.

\begin{table}[H]
\centering
\label{table:all-scenarios:start}
\begin{tabular}{lllrrrr}
\toprule
\textbf{Arrival} & \textbf{Network} & \textbf{Scheduler} & \textbf{A01} & \textbf{A02} & \textbf{A03} & \textbf{On time} \\
\midrule
Steady & Normal & Round Robin & 181.7 & 236.7 & 227.2 & 36/36 \\
Steady & Normal & Compute-ECT & 118.0 & 121.2 & 119.2 & 36/36 \\
Steady & Normal & Network-ECT & 113.6 & 120.1 & 119.2 & 36/36 \\
Steady & +20 ms & Round Robin & 266.1 & 252.4 & 248.9 & 36/36 \\
Steady & +20 ms & Compute-ECT & 209.0 & 203.5 & 199.1 & 36/36 \\
Steady & +20 ms & Network-ECT & 196.3 & 207.0 & 145.7 & 36/36 \\
Burst & Normal & Round Robin & 185.5 & 171.9 & 189.8 & 35/36 \\
Burst & Normal & Compute-ECT & 131.2 & 143.3 & 136.0 & 36/36 \\
Burst & Normal & Network-ECT & 137.3 & 138.3 & 139.0 & 36/36 \\
Burst & +20 ms & Round Robin & 210.6 & 210.4 & 210.2 & 36/36 \\
Burst & +20 ms & Compute-ECT & 191.0 & 186.8 & 198.7 & 36/36 \\
Burst & +20 ms & Network-ECT & 186.9 & 206.5 & 204.5 & 36/36 \\
\bottomrule
\end{tabular}
\label{table:all-scenarios:end}
\end{table}


\begin{figure}[H]
\centering
\includegraphics[width=1\linewidth,height=\textheight,keepaspectratio,alt={Median latency in all three VM sessions; the connecting range is descriptive, not a confidence interval}]{build/figures/median-latency.pdf}
\caption{Median latency in all three VM sessions; the connecting range
is descriptive, not a confidence interval}
\end{figure}

\subsection{5. What did I measure?}\label{what-did-i-measure}

The procedure and dependency baseline were specified before collection
and independently reviewed. Each VM used the same pinned software,
four-container topology, CPU limits, task streams and scenario order.
\textbf{Each session was attempted once, without retries or tuning.} The
public repository provides the execution source, package pins and
commands.

Before application work, I check the emulator\textquotesingle s delay,
bandwidth limit, routes, disconnection/reconnection and unaffected path.
Application correctness and network fidelity are assessed separately.
After collection, I verify the evidence export and VM cleanup before
finalizing eligibility for comparison.

One full session contains:

\begin{table}[H]
\centering
\label{table:stages:start}
\begin{tabular}{lrl}
\toprule
\textbf{Stage} & \textbf{Per VM} & \textbf{What it checks} \\
\midrule
Calibration & 72 & Estimate worker and transfer times \\
Queue check & 12 & Confirm waiting behind busy workers \\
Warmup & 720 & Prepare every scenario equally \\
Main comparison & 144 & 12 scenarios with 12 scored tasks each \\
Total real work & 948 & Includes preparation and scored work \\
\bottomrule
\end{tabular}
\label{table:stages:end}
\end{table}


I record correct results, deadline success, end-to-end latency, worker
selection, queue waiting and scheduling overhead. \textbf{Decision wall
time} is elapsed time spent deciding; \textbf{decision CPU time} counts
CPU used by that thread. Worker execution and communication time remain
separate.

\begin{figure}[H]
\centering
\includegraphics[width=1\linewidth,height=\textheight,keepaspectratio,alt={Scheduling decision wall time in all three sessions; each point is one scenario median}]{build/figures/decision-overhead.pdf}
\caption{Scheduling decision wall time in all three sessions; each point
is one scenario median}
\end{figure}

\Needspace{100pt}

Across the 36 scenario/session points, decision wall medians ranged from
\textbf{9.375 to 160.792 microseconds}; decision CPU medians ranged from
\textbf{6.459 to 146.083 microseconds}. In all 24 matched
ECT-minus-Round-Robin comparisons, the ECT decision wall and CPU medians
were higher. This overhead is retained alongside the lower response
medians; it is not subtracted from the measured response.

\subsection{6. What happened in the
experiment?}\label{what-happened-in-the-experiment}

All three prospectively listed sessions completed their 12 scenarios.
All scored tasks returned correct answers; none was rejected, failed or
unresolved.

\begin{table}[H]
\centering
\label{table:sessions:start}
\begin{tabular}{llrr}
\toprule
\textbf{Session} & \textbf{Application} & \textbf{Correct outputs} & \textbf{Scored on time} \\
\midrule
Run 1 & Complete & 948 & 144/144 \\
Run 2 & Complete & 948 & 144/144 \\
Run 3 & Complete & 948 & 143/144 \\
\bottomrule
\end{tabular}
\label{table:sessions:end}
\end{table}


The Gateway recorded \textbf{948 assignments per session}, including 144
scored decisions. The worker totals below include calibration, queue
checks and warmup; scored counts are shown in parentheses.

\begin{table}[H]
\centering
\label{table:nodes:start}
\begin{tabular}{lrrrr}
\toprule
\textbf{Session} & \textbf{G all / scored} & \textbf{Edge} & \textbf{Fog} & \textbf{Cloud} \\
\midrule
Run 1 & 948 / 144 & 284 (16) & 294 (26) & 370 (102) \\
Run 2 & 948 / 144 & 284 (16) & 294 (26) & 370 (102) \\
Run 3 & 948 / 144 & 284 (16) & 298 (30) & 366 (98) \\
\bottomrule
\end{tabular}
\label{table:nodes:end}
\end{table}


Both health and application network checks passed the original limits in
all three sessions. The application-quota measurements were:

\begin{table}[H]
\centering
\label{table:network:start}
\begin{tabular}{lrrr}
\toprule
\textbf{Session} & \textbf{Forward delta (ms)} & \textbf{Reverse delta (ms)} & \textbf{Limited rate (Mbps)} \\
\midrule
Run 1 & 22.139 & 20.779 & 9.072 \\
Run 2 & 21.223 & 21.570 & 9.320 \\
Run 3 & 22.045 & 21.967 & 8.357 \\
\bottomrule
\end{tabular}
\label{table:network:end}
\end{table}


The nominal settings were 20 ms added delay and a separate 10 Mbps
qualification limit. The original accepted bands were \textbf{15-25 ms}
for the delay increment and \textbf{7-11 Mbps} for the limited receive
rate; baseline, restoration, route and unaffected-path checks also
passed. The rate limit was a network qualification check, not a third
scheduling condition.

\Needspace{120pt}

\textbf{Repeatability has limits.} Across the three VM medians for a
fixed scenario, the relative range was \textbf{0.219\%-31.267\%},
calculated as (maximum minus minimum) divided by their median. The
largest spread was Network-ECT with steady arrivals and added Cloud
delay: \textbf{196.273, 207.037 and 145.668 ms}. The low A03 value
remains visible in the graph. Three complete sessions establish this
bounded correctness result, not identical timing or general future
reliability.

All scenarios completed 12 tasks inside the six-second window, giving
\textbf{2 completed tasks/s}. This small offered workload does not
measure maximum processing capacity. On-time goodput was 2 tasks/s
except the A03 late-task scenario, where it was 11/6 tasks/s.

\subsection{7. How do individual tasks
behave?}\label{how-do-individual-tasks-behave}

A \emph{timeline} shows when a task was released, waited, transferred,
executed and returned. The figure uses \textbf{all 12 actual tasks from
A03/Round Robin/burst/normal network}, chosen to show the only deadline
miss rather than hide it. Each colored segment is measured elapsed time;
the small scheduling stages remain included.

\begin{figure}[H]
\centering
\includegraphics[width=1\linewidth,height=\textheight,keepaspectratio,alt={A03 task timeline for Round Robin with burst arrivals; the late task is retained}]{build/figures/task-timeline.pdf}
\caption{A03 task timeline for Round Robin with burst arrivals; the late
task is retained}
\end{figure}

\textbf{The late Edge task} took \textbf{1,002.318 ms} from planned
arrival to verified response, exceeding the 1,000 ms deadline by
\textbf{2.318 ms}. Its actual release was 4.019 ms late. Measuring from
that actual release would give 998.299 ms and hide the miss, so the
original planned-arrival rule is retained. Its answer was correct.

The measurements do not isolate the cause of this delay. Worker CPU
limits, host/guest scheduling and measurement activity can all affect
elapsed time; this study did not vary them independently.

\Needspace{260pt}

\subsection{8. Current conclusion}\label{current-conclusion}

\textbf{Demonstrated:} the same application orchestrator completed in
\textbf{three fresh VM sessions on one host}, with correct task
accounting and outputs on every worker, separately qualified network
conditions, intact evidence and verified cleanup. Across the full
procedure there were \textbf{2,844 correct worker results}; among the
\textbf{432 scored offers, 431 met the deadline and one was late}.

\textbf{Still limited:} this is three small sessions on one physical
host, with a fixed task stream and scenario order. Timing and placements
are not identical. Fog couples routing and computation. Emulated links
and CPU quotas are not separate physical sites. These data do not
establish a scheduler winner, geographic portability or long-term
operational reliability.

\textbf{Next step:} use this reproducible prototype to agree a focused
question with the supervisors, then specify a broader workload or
validation setting before collecting more data. Increasing load,
changing network information or adding physical hosts would be separate
experimental decisions.

\emph{This report uses only the recorded A01-A03 cohort. The figures and
tables were generated offline; no new experiment was run to prepare this
update.}

\subsection{References}\label{references}

\begin{enumerate}
\def\labelenumi{\arabic{enumi}.}
\tightlist
\item
  Z. Wang, B. Sedlak, J. L. Herrera, and S. Dustdar.
  \textquotesingle Fair Comparison of Scheduling Algorithms on
  Heterogeneous Edge Clusters: A Continuous Adaptive
  Benchmark.\textquotesingle{} IEEE ICDCSW, pp. 232-237, 2026.
  \href{https://doi.org/10.1109/ICDCSW72724.2026.00054}{DOI:
  10.1109/ICDCSW72724.2026.00054}. Open preprint:
  \href{https://arxiv.org/abs/2606.12343}{arXiv:2606.12343}.
\item
  MADT project. \textquotesingle Distributed Application Modelling
  System.\textquotesingle{}
  \href{https://github.com/dltcspbu/madt}{Source repository}, accessed
  23 October 2026.
\item
  \href{https://github.com/HronoSF/mozhaiskii-edge-fog-cloud-research}{Public
  experiment source, data and execution instructions}. A01-A03 data
  release, October 2026.
\end{enumerate}

\end{document}
